\documentclass[a4paper,12pt]{article}

\usepackage[T2A]{fontenc}
\usepackage[utf8]{inputenc}
\usepackage[russian]{babel}
\usepackage{amsmath,amssymb,mathtools}
\usepackage{PTSans}
\renewcommand{\familydefault}{\sfdefault}
\usepackage{icomma}
\usepackage{geometry}
\usepackage{microtype}
\usepackage{tikz}
\usetikzlibrary{calc}
\usepackage{tabularx,booktabs,longtable}
\usepackage{enumitem}
\usepackage{hyperref}
\usepackage{parskip}
\usepackage{fancyhdr}
\usepackage{setspace}
\usepackage{needspace}
\usepackage{xcolor}

\geometry{left=20mm,right=20mm,top=18mm,bottom=20mm}
\setlength{\headheight}{25pt}
\onehalfspacing
\setlength{\parindent}{0pt}
\setlist[itemize]{leftmargin=6mm,itemsep=2pt,topsep=3pt}
\setlist[enumerate]{leftmargin=8mm,itemsep=5pt,topsep=4pt}

\definecolor{accent}{HTML}{147D7E}
\definecolor{darkaccent}{HTML}{0B4F50}
\definecolor{textgray}{HTML}{303638}
\definecolor{warn}{HTML}{8B5E3C}

\hypersetup{colorlinks=true,linkcolor=darkaccent,urlcolor=darkaccent}

\pagestyle{fancy}
\fancyhf{}
\fancyhead[L]{\small\color{textgray}Повторение ключевых тем}
\fancyhead[R]{\small\color{textgray}05.10.26}
\fancyfoot[C]{\small\color{textgray}\thepage}
\renewcommand{\headrulewidth}{0.3pt}
\renewcommand{\headrule}{\hbox to\headwidth{\color{accent}\leaders\hrule height \headrulewidth\hfill}}

\newcommand{\topic}[1]{%
  \Needspace{8\baselineskip}%
  \vspace{0.6em}
  {\Large\bfseries\color{darkaccent}#1}\par
  \vspace{0.25em}
  {\color{accent}\rule{\linewidth}{0.6pt}}\par
  \vspace{0.3em}
}
\newcommand{\key}[1]{\textbf{\color{darkaccent}#1}}
\newcommand{\mistake}[1]{\textbf{\color{warn}Ошибка:} #1}
\newcommand{\checkitem}{\(\square\)}

\renewcommand{\today}{05.10.26}

\begin{document}

\thispagestyle{empty}
\vspace*{24mm}
\begin{center}
  {\fontsize{30}{34}\selectfont\bfseries\color{darkaccent} Чек-лист}\\[4mm]
  {\fontsize{21}{25}\selectfont\bfseries Повторение ключевых тем ЕГЭ}\\[8mm]
  {\Large Геометрия, векторы, вероятность, уравнения, текстовые задачи}\\[15mm]

  {\large Ксения Клыкова}\\[3mm]
  {\large \today}\\[18mm]

  \begin{minipage}{0.78\textwidth}
  \centering
  \color{textgray}
  Цель: восстановить формулы и алгоритмы, которые уже встречались,
  и научиться быстро узнавать тип задачи по условию.
  \end{minipage}
\end{center}

\vfill

\begin{center}
{\large\bfseries\color{darkaccent}Лёвин Артём Александрович эксклюзивно для Ксении Клыковой}
\end{center}

\vspace{8mm}
\begin{tikzpicture}[remember picture,overlay]
  \draw[accent,line width=1.2pt]
    ([xshift=18mm,yshift=17mm]current page.south west)
    .. controls ([xshift=65mm,yshift=29mm]current page.south west)
    and ([xshift=-60mm,yshift=7mm]current page.south east)
    .. ([xshift=-18mm,yshift=20mm]current page.south east);
\end{tikzpicture}

\clearpage

\topic{1. Окружность и углы}

\key{Вписанный угол.}
Если вершина угла лежит на окружности, то его градусная мера равна половине градусной меры дуги, на которую он опирается:
\[
\angle ADB=\frac12\,\overset{\frown}{AB}.
\]

\key{Сумма углов треугольника.}
После нахождения нужного угла через дугу часто остаётся обычный треугольник:
\[
\angle A+\angle C+\angle D=180^\circ.
\]

\begin{itemize}
  \item \checkitem{} Сначала определите, на какую именно дугу опирается вписанный угол.
  \item \checkitem{} Убедитесь, что выбрана дуга, не содержащая вершину угла.
  \item \checkitem{} После перехода к треугольнику используйте сумму углов.
\end{itemize}

\mistake{увидеть подписанный угол или дугу, но не включить их в решение. Перед вычислениями перенесите все числовые данные на схему.}

\topic{2. Площадь треугольника через две стороны и угол}

Если известны две стороны \(a,b\) и угол \(C\) между ними:
\[
S=\frac12 ab\sin C.
\]

Для равнобедренного треугольника с боковыми сторонами \(x\) и углом между ними \(30^\circ\):
\[
S=\frac12 x^2\sin30^\circ=\frac{x^2}{4}.
\]

Полезные значения:
\[
\sin30^\circ=\frac12,\qquad
\cos45^\circ=\frac{\sqrt2}{2}.
\]

\begin{itemize}
  \item \checkitem{} Проверьте, что используемый угол действительно расположен между выбранными сторонами.
  \item \checkitem{} Если после формулы появляется \(x^2=c\), длина стороны берётся положительной.
\end{itemize}

\topic{3. Скалярное произведение векторов}

Основная формула:
\[
\vec a\cdot\vec b=|\vec a|\,|\vec b|\cos\varphi,
\]
где \(\varphi\) --- угол между векторами.

Если требуется длина второго вектора:
\[
|\vec b|=\frac{\vec a\cdot\vec b}{|\vec a|\cos\varphi}.
\]

\begin{itemize}
  \item \checkitem{} Выпишите формулу до подстановки чисел.
  \item \checkitem{} Вспомните точное значение тригонометрической функции.
  \item \checkitem{} Длина вектора не может быть отрицательной.
\end{itemize}

\topic{4. Вероятность: базовая модель}

Для равновозможных исходов:
\[
P(A)=\frac{\text{число благоприятных исходов}}{\text{общее число исходов}}.
\]

Обязательная самопроверка:
\[
0\le P(A)\le1.
\]

Если получилось \(1{,}2\), \(-0{,}3\) или любое другое значение вне отрезка \([0;1]\),
арифметику или модель события нужно перепроверить.

\topic{5. Дерево вероятностей и условная вероятность}

На дереве вероятностей:

\begin{itemize}
  \item вдоль одной траектории вероятности \key{перемножаются};
  \item вероятности нескольких подходящих траекторий \key{складываются}.
\end{itemize}

Условная вероятность:
\[
P(A\mid B)=\frac{P(A\cap B)}{P(B)},\qquad P(B)>0.
\]

Алгоритм:
\begin{enumerate}
  \item Найдите все траектории, соответствующие событию \(B\), и вычислите \(P(B)\).
  \item Найдите траектории, где одновременно произошли \(A\) и \(B\), и вычислите \(P(A\cap B)\).
  \item Разделите \(P(A\cap B)\) на \(P(B)\).
  \item Проверьте, что ответ лежит в \([0;1]\).
\end{enumerate}

\mistake{делить на \(P(A)\) только потому, что событие \(A\) записано первым. В знаменателе стоит вероятность условия --- события после вертикальной черты.}

\topic{6. Формула Бернулли}

Вероятность получить ровно \(m\) успехов в \(n\) независимых испытаниях:
\[
P_n(m)=\binom{n}{m}p^m(1-p)^{n-m},
\qquad
\binom{n}{m}=\frac{n!}{m!(n-m)!}.
\]

Для честной монеты \(p=\frac12\), поэтому:
\[
P_n(m)=\binom{n}{m}\left(\frac12\right)^n.
\]

Если спрашивают, \key{во сколько раз} одна вероятность больше другой, удобно сразу составить отношение. Общие множители сокращаются.

При сокращении факториалов раскрывайте только нужную часть:
\[
10!=10\cdot9\cdot8\cdot7\cdot6\cdot5!.
\]

\mistake{раскрывать все факториалы до единицы и получать громоздкие числа. Сначала ищите общие множители для сокращения.}

\clearpage

\topic{7. Уравнения с корнем}

Для уравнения вида
\[
\sqrt{x-a}=b
\]
сначала требуется \(b\ge0\), затем:
\[
x-a=b^2.
\]

После возведения в квадрат полезно подставить найденное значение в исходное уравнение.

Пример структуры:
\[
\sqrt{x-2}=6
\quad\Longrightarrow\quad
x-2=36
\quad\Longrightarrow\quad
x=38.
\]

\topic{8. Равенство дробей и ОДЗ}

Если
\[
\frac{1}{A(x)}=\frac{1}{B(x)},
\]
то при условиях
\[
A(x)\ne0,\qquad B(x)\ne0
\]
можно сделать вывод:
\[
A(x)=B(x).
\]

Сначала зафиксируйте ограничения, затем решайте полученное линейное уравнение.

\mistake{забыть проверить, что найденный корень не обращает знаменатель в ноль.}

\topic{9. Прикладные задачи по формуле}

Рабочий порядок:
\begin{enumerate}
  \item Внимательно прочитайте, какую величину требуется найти.
  \item Выпишите формулу.
  \item Подставьте данные вместе с единицами измерения, если они есть.
  \item Выполните вычисления по шагам.
  \item Если получили, например, \(\sqrt{0{,}09}\), явно запишите переход:
  \[
  \sqrt{0{,}09}=0{,}3.
  \]
  \item Сверьте смысл и размер ответа с условием.
\end{enumerate}

\topic{10. Текстовые задачи на движение}

Главный инструмент --- таблица
\[
\begin{array}{c|c|c|c}
 & S & v & t\\
\hline
\text{1-й объект} & \cdots & \cdots & \cdots\\
\text{2-й объект} & \cdots & \cdots & \cdots
\end{array}
\]
и связь
\[
S=vt,\qquad t=\frac{S}{v}.
\]

Если первый автомобиль выехал раньше на \(3\) часа, то к моменту встречи он двигался на \(3\) часа дольше:
\[
t_1-t_2=3.
\]

Если \(x\) --- расстояние от города \(A\) до места встречи, а всё расстояние равно \(510\) км:
\[
S_1=x,\qquad S_2=510-x.
\]

\mistake{перепутать, чьё время больше. Раньше выехал \(\Rightarrow\) дольше ехал.}

\clearpage

\section*{\color{darkaccent}Тренировочный блок --- Путь А}

\textbf{10 заданий.} Решения оформляйте кратко, но так, чтобы был виден использованный метод.

\begin{enumerate}
  \item Дуга \(AB\) окружности равна \(124^\circ\). Найдите вписанный угол \(ADB\), опирающийся на эту дугу.

  \item В треугольнике \(ACD\) известно:
  \(\angle ACD=42^\circ\), \(\angle CDA=62^\circ\).
  Найдите \(\angle CAD\).

  \item В равнобедренном треугольнике угол между боковыми сторонами равен \(30^\circ\), а площадь равна \(25\).
  Найдите длину боковой стороны.

  \item Известно:
  \[
  |\vec a|=3,\qquad
  \angle(\vec a,\vec b)=45^\circ,\qquad
  \vec a\cdot\vec b=\frac{15\sqrt2}{2}.
  \]
  Найдите \(|\vec b|\).

  \item В коробке лежат \(7\) красных, \(5\) синих и \(3\) зелёных шара.
  Наугад выбирают один шар. Найдите вероятность того, что он зелёный.

  \item Событие \(A\) происходит с вероятностью \(0{,}6\).
  Если произошло \(A\), событие \(B\) происходит с вероятностью \(0{,}4\);
  если \(A\) не произошло, событие \(B\) происходит с вероятностью \(0{,}8\).
  Найдите \(P(A\mid B)\).

  \item Честную монету бросают \(10\) раз.
  Во сколько раз вероятность получить ровно \(5\) орлов больше вероятности получить ровно \(4\) орла?

  \item Решите уравнение:
  \[
  \sqrt{x-2}=6.
  \]

  \item Решите уравнение:
  \[
  \frac{1}{7x+16}=\frac{1}{8x+11}.
  \]

  \item Расстояние между городами \(A\) и \(B\) равно \(510\) км.
  Из \(A\) в \(B\) выехал автомобиль со скоростью \(70\) км/ч.
  Через \(3\) часа навстречу ему из \(B\) выехал второй автомобиль со скоростью \(80\) км/ч.
  На каком расстоянии от города \(A\) автомобили встретятся?
\end{enumerate}

\clearpage

\section*{\color{darkaccent}Ответы}

\renewcommand{\arraystretch}{1.35}
\begin{tabularx}{\linewidth}{@{}>{\bfseries}c X@{}}
\toprule
№ & Ответ \\
\midrule
1 & \(62^\circ\). \\
2 & \(76^\circ\). \\
3 & \(10\). \\
4 & \(5\). \\
5 & \(\dfrac15=0{,}2\). \\
6 & \(\dfrac37\approx0{,}429\). \\
7 & \(\dfrac65=1{,}2\) раза. \\
8 & \(38\). \\
9 & \(5\). \\
10 & \(350\) км. \\
\bottomrule
\end{tabularx}

\vspace{10mm}

{\small\color{textgray}
Контроль перед завершением: вероятность находится в диапазоне \([0;1]\);
знаменатели не равны нулю; длины положительны; в задаче на движение большее время относится к автомобилю, который стартовал раньше.
}

\end{document}
